% Authoritative LaTeX manuscript; Markdown is a generated compatibility projection.

A working control becomes part of the system later changes must preserve. Its inputs can drift, its consumer can move, and another control can interfere with it. The preceding blueprints therefore need maintenance beyond keeping their files present: engineers must examine interactions, correspondence, failure responses, and the work the machinery actually performs.

\section{Analyze Interactions Before Adding Another Control}\label{heading:handbook-model-governance-graph}

Adding a control can create interactions with existing machinery. Two hooks may write the same record; one mechanism may need to inspect state before another changes it; or a required property of a commit may conflict with another transformation. The useful model is therefore not a catalogue of controls but a representation of the relationships needed to decide whether they interfere.

% Evidence: E24 | source: book/data/local-evidence/evidentiary-basis-11202026/integration/records/E24.md | kind: implementation
DocAble's governance graph represents selected mechanisms with firing events, resource requirements, enforcement characteristics, and implementation anchors. Requirements distinguish reading, writing, requiring a property, and constraining a resource. Resources have hierarchical names: a coarse reference to \texttt{commit} can be refined to \texttt{commit.message} or \texttt{commit.body}. Apparent interference may disappear when the modeled subresources are disjoint. Genuine overlap may require an ordering or an engineering decision about incompatible requirements.

% Evidence: E24 | source: book/data/local-evidence/evidentiary-basis-11202026/integration/records/E24.md | kind: implementation
The graph's analysis computes a declared order and detects cycles; tests include both a two-node cycle and a self-loop. This establishes properties of the represented relationships. A query returning an order is not an executing scheduler. If an analogous arrangement requires enforced ordering, install a consumer that actually uses the result and test its observed execution sequence.

% Evidence: E24 | source: book/data/local-evidence/evidentiary-basis-11202026/integration/records/E24.md | kind: implementation
Correspondence has a similarly concrete boundary. The anchor check compares modeled commands, lifecycle events, scripts, and entry symbols with configured wiring. It can detect a missing command or a mismatched firing event within its inspected population. Finding the expected function name does not establish that its behavior remains suitable. When a control's meaning changes, review the represented requirements as well as its address.

\section{Reuse Machinery Across Obligations}\label{heading:handbook-reusable-evaluation}

% Evidence: AA1 | source: book/data/local-evidence/automated-alignment-11242026/records/AA1.md | kind: implementation
A recurring class of obligations can justify a common representation and evaluator. DocAble's component model declares permitted imports for selected Python components; one checker applies those declarations to the corresponding source files. An additional boundary extends the checked population without another checker. Construct such a mechanism around explicit identities, the meaning of each declaration, and a rule for selecting its population. Test both a forbidden import and an allowed one, and distinguish a component with no declared boundary from one whose imports satisfy its boundary.

% Evidence: AA2 | source: book/data/local-evidence/automated-alignment-11242026/records/AA2.md | kind: implementation
Models can also determine what verification machinery must exist. DocAble derives a verification tier from an invariant's participating lanes, coordination primitive, and temporal operator. Its policy assigns concurrent safety obligations to explicit-state exploration, concurrent liveness obligations to a temporal specification, and ordinary sequential obligations to property tests. A reusable lint maps each tier to a required artifact kind and checks that a reference of that kind resolves to a file. The mapping must cover every tier; a new tier without a mapping is itself a finding. The machinery applies across models, with blocking enforcement for selected models and advisory findings while others are brought into conformance.

This construction separates three decisions: how to classify an obligation, which artifact its class requires, and whether that artifact supplies adequate evidence. The first two can follow an explicit policy; the last still requires examination of the test or specification. Challenge the mechanism with a missing reference, a wrong artifact kind, a nonexistent file, and an unmapped tier. A vacuous checker can pass an existence check, so artifact presence must not stand for adequate verification.

% Evidence: AA2 | source: book/data/local-evidence/automated-alignment-11242026/records/AA2.md | kind: implementation
Keep artifact correspondence separate from execution as well. DocAble derives a temporal property definition from its model and checks that the TLA+ specification contains the corresponding property; the state and action definitions require separate construction. The parity check is blocking, but running TLC remains a separate activity. If an adoption requires analysis on every relevant change, it also needs an execution gate and evidence that the analysis runs on the candidate being admitted.

\section{Recurring Engineering Moves}\label{heading:handbook-composition-procedure}

% Evidence: E10 | source: book/data/local-evidence/evidentiary-basis-11202026/integration/records/E10.md | kind: implementation
The ten moves in \autoref{handbook-engineering-moves} are construction choices. A blueprint combines the moves needed for its purpose; it does not require every mechanism in the table. For example, deriving a grant plan is useful only after the caller, callee, and runtime identity can be joined reliably. Checking that derivation against an existing deployment map is a migration aid, not another permanent owner of the same fact.

{\footnotesize
\begin{longtable}{@{}p{0.24\linewidth}p{0.28\linewidth}p{0.38\linewidth}@{}}
\caption[Engineering moves at a glance]{Ten engineering moves, the problems that motivate them, and concrete realizations in DocAble.}\label{handbook-engineering-moves}\tabularnewline
\toprule
Recurring problem & Engineering move & DocAble example \\
\midrule
\endfirsthead
\toprule
Recurring problem & Engineering move & DocAble example \\
\midrule
\endhead
% Evidence: E19 | source: book/data/local-evidence/evidentiary-basis-11202026/integration/records/E19.md | kind: implementation
% Evidence: E34 | source: book/data/local-evidence/evidentiary-basis-11202026/integration/records/E34.md | kind: implementation
Different consumers disagree about the same fact. & Make one source authoritative. Assign the fact an owner and name its consumers. & The component catalogue defines the zones used by repository accounting and agent audits. \\
% Evidence: E10 | source: book/data/local-evidence/evidentiary-basis-11202026/integration/records/E10.md | kind: implementation
A copied declaration drifts from its source. & Derive; do not copy. Compute the second representation where a reliable transformation exists. & Deployment derives invocation-grant plans from service dependencies and runtime identity. \\
% Evidence: E24 | source: book/data/local-evidence/evidentiary-basis-11202026/integration/records/E24.md | kind: implementation
A useful representation no longer describes its subject. & Check correspondence. Compare declared facts with independently obtained evidence. & A governance-graph check resolves modeled hooks against their configured events and implementing scripts. \\
% Evidence: E36 | source: book/data/local-evidence/evidentiary-basis-11202026/integration/records/E36.md | kind: implementation
A collection of tests has no explicit account of what it must cover. & Derive the obligation set. Enumerate required evidence from represented properties. & Cross-service invariant declarations identify verification obligations and references to their checkers. \\
% Evidence: E09 | source: book/data/local-evidence/evidentiary-basis-11202026/integration/records/E09.md | kind: implementation
% Evidence: E15 | source: book/data/local-evidence/evidentiary-basis-11202026/integration/records/E15.md | kind: implementation
A rule is checked where the necessary state is unavailable. & Enforce where the property is legible. Place the decision and its consequence at a suitable boundary. & The PDF applier resolves the target before deciding whether an artifact-demotion request is admissible. \\
% Evidence: E15 | source: book/data/local-evidence/evidentiary-basis-11202026/integration/records/E15.md | kind: implementation
A protected operation remains available through ungoverned paths. & Close the action surface. Offer supported operations through a controlled interface. & External PDF editing uses typed operations; the applier can refuse a request before mutation. \\
% Evidence: E19 | source: book/data/local-evidence/evidentiary-basis-11202026/integration/records/E19.md | kind: implementation
% Evidence: E58 | source: book/data/local-evidence/evidentiary-basis-11202026/integration/records/E58.md | kind: implementation
Relevant knowledge is discovered after the consequential choice. & Deliver knowledge at the decision. Join the work's scope to the representations it needs. & Dispatch context extraction selects constraints associated with the brief's target files. \\
% Evidence: E14 | source: book/data/local-evidence/evidentiary-basis-11202026/integration/records/E14.md | kind: implementation
A result cannot be traced to the work that produced it. & Carry cause with consequence. Preserve identities and attribution across transformations. & PDF mutation records identify the responsible pass or editing session; applied results can be compared with proposed edits. \\
% Evidence: E32 | source: book/data/local-evidence/evidentiary-basis-11202026/integration/records/E32.md | kind: practice-and-design
% Evidence: E60 | source: book/data/local-evidence/evidentiary-basis-11202026/integration/records/E60.md | kind: practice-and-design
The same operational judgment must be reconstructed repeatedly. & Externalize recurring judgment. Record states, evidence, and sanctioned responses. & Lifecycle descriptions and runbooks distinguish recovery work from ordinary completion and identify where judgment remains. \\
% Evidence: E24 | source: book/data/local-evidence/evidentiary-basis-11202026/integration/records/E24.md | kind: implementation
Individually reasonable changes conflict when combined. & Make dependencies queryable. Represent the relationships needed to analyze the interaction. & The governance graph records firing events, resource use, and chosen resolutions among selected mechanisms. \\
\bottomrule
\end{longtable}
}

\section{Repair, Reconsider, or Retire}

% Evidence: E14 | source: book/data/local-evidence/evidentiary-basis-11202026/integration/records/E14.md | kind: implementation
% Evidence: E57 | source: book/data/local-evidence/evidentiary-basis-11202026/integration/records/E57.md | kind: incident
A control's repeated use supplies evidence about its design. False refusals may reveal a bad identity extractor, as in the completed-phase gate. Missing attribution may reveal a mutation path the scanner cannot enumerate. Repeated manual reconstruction may identify a useful model; a frequently consulted model that never changes a decision may instead be unnecessary. Distinguish a defect in the mechanism from a deliberate exception to its policy before adding another override or check.

When a control needs repair, identify the failure class, the mechanism responsible, and the evidence needed to test the correction. Retest both prohibited and legitimate work. Retire superseded declarations and controls once their consumers have migrated; leaving both active can create conflicting authority, duplicated cost, or misleading evidence of independent protection.
